

\section{Experiments} \label{sec:exp}
\subsection{Experimental Setups}
\noindent\textbf{Benchmarks.}
We evaluate routers across the five xRouteBench tracks defined in \S\ref{sec:xroutebench}: Generic LLM Tasks, memory, vision, time-series, and personalized. Together, they comprise eight test sets, with full statistics reported in Table~\ref{tab:app-benchmarks}. For each query in the memory track, we retrieve up to five memory items as context and score responses using token-level F1.

\noindent\textbf{LLM Candidates.}
The candidate pool contains 18 open-weight models served through two providers (i.e., Together API\footnote{\url{https://www.together.ai/}} and NVIDIA NIM API\footnote{\url{https://build.nvidia.com/}}), spanning 7B to 671B parameters. It covers Gemma-2-9B \citep{team2024gemma}; Mistral-7B, Mistral-Small-24B, Mixtral-8x7B, and Mixtral-8x22B \citep{jiang2023mistral, jiang2024mixtral}; Qwen2.5-7B, Qwen3-Next-80B, and Qwen3-Coder \citep{yang2024qwen2, yang2025qwen3}; Llama-3-8B, Llama-3-70B, Llama-3.3-70B, and Llama-4-Maverick \citep{grattafiori2024llama, adcock2026llama}; GPT-OSS-20B and GPT-OSS-120B \citep{agarwal2025gpt}; RNJ-1-15B \citep{rnj1_5_instruct}; and the two 671B models DeepSeek-V3.1 \citep{liu2024deepseek} and Cogito-v2 \citep{deepcogito2025cogitov2_671b}. Per-token pricing is given in Table~\ref{tab:app-pricing} in the appendix.

\noindent\textbf{Implemented Routers.}
\method\ implements more than 16 routers covering three families. (i) \textbf{Single-turn routers} include $k$NNRouter~\citep{li2025rethinking}, SVMRouter, MLPRouter, EloRouter, and MFRouter~\citep{ong2024routellm, shnitzer2023large}, as well as RouterDC \citep{chen2024routerdc}, Hybrid LLM \citep{ding2024hybrid}, AutoMix \citep{aggarwal2024automix}, GraphRouter \citep{feng2024graphrouter}, CausalLM~\citep{ong2024routellm}, and two rule-based baselines that always select the smallest or largest model; (ii) \textbf{Multi-turn routers} include Router-R1 \citep{zhang2025router} together with $k$NN-based and LLM-based multi-round routers; and (iii) \textbf{Personalized routers} include GMTRouter~\citep{xie2025gmtrouter} and PersonalizedRouter~\citep{dai2025personalizedrouter}.

\noindent\textbf{Evaluation Protocol.}
We score each router by a weighted reward $\alpha \cdot \mathrm{perf} - \beta \cdot \mathrm{cost}$. We sweep five weight settings from the quality-only $(\alpha, \beta) = (1.0, 0.0)$ to the heavily cost-weighted $(0.2, 0.8)$. Multi-round and RL-based routers cannot optimize this weighted objective and are therefore run once under a single configuration. On the personalized track, answers are scored by a persona-conditioned LLM judge (DeepSeek-V3.1) as win, tie, or loss ($1$, $0.5$, $0$), and the judge's own cost is excluded from the reported cost. 

\subsection{Main Results}
\label{sec:results}

\begin{table}[t]
\centering
\caption{\textbf{Results on xRouteBench under the performance-first setting $(\alpha, \beta) = (1.0, 0.0)$.} Scores are reported across the Generic LLM Tasks, memory, vision, and time-series
tracks, together with their average. Following the original implementations where applicable, all multi-turn routers use Qwen2.5-3B-Instruct as the base model. Top two results are highlighted in \textbf{bold} and \underline{underline}.}
\label{tab:main}
\small
\resizebox{\linewidth}{!}{%
\begin{tabular}{lcccccccc}
\toprule
\multirow{2}{*}{\textbf{Router}} & \multirow{2}{*}{\textbf{Generic LLM Tasks}} & \multicolumn{2}{c}{\textbf{Memory}} & \multicolumn{3}{c}{\textbf{Vision}} & \multirow{2}{*}{\textbf{TimeSeries}} & \multirow{2}{*}{\textbf{Avg}} \\
\cmidrule(lr){3-4}\cmidrule(lr){5-7}
 & & \textbf{LoCoMo} & \textbf{LongMemEval} & \textbf{Geometry3K} & \textbf{MathVista} & \textbf{Video} & & \\
\midrule
\rowcolor{barGray} \multicolumn{9}{c}{\textit{Rule-based baselines}} \\
Smallest-LLM & 57.55 & 25.44 & 36.77 & 27.87 & 35.00 & \textbf{33.33} & 49.61 & 37.94 \\
Largest-LLM & 70.29 & 26.59 & 35.57 & 37.70 & 33.00 & 22.22 & 45.67 & 38.72 \\
\midrule
\rowcolor{barGray} \multicolumn{9}{c}{\textit{Single-turn routers}} \\
$k$NNRouter & 71.37 & 25.24 & \textbf{38.74} & 31.15 & 41.00 & \underline{29.63} & 51.97 & 41.30 \\
SVMRouter & 74.21 & \textbf{27.64} & \underline{38.68} & \underline{42.62} & \underline{47.00} & \underline{29.63} & 55.91 & \underline{45.10} \\
MLPRouter & 68.12 & \underline{26.78} & 32.27 & 27.87 & 34.00 & \underline{29.63} & 56.69 & 39.34 \\
MFRouter & 67.23 & 24.49 & 34.91 & 40.98 & 29.00 & 22.22 & 51.97 & 38.69 \\
EloRouter & 64.15 & 25.70 & 37.27 & \textbf{45.90} & \textbf{50.00} & 25.93 & \textbf{63.78} & 44.68 \\
Hybrid LLM & 64.68 & 25.89 & 36.56 & 32.79 & 37.00 & \textbf{33.33} & 51.18 & 40.20 \\
RouterDC & \textbf{80.56} & 24.93 & 36.77 & 16.39 & 24.00 & 25.93 & 45.67 & 36.32 \\
GraphRouter & \underline{80.54} & 25.94 & 33.93 & \underline{42.62} & \textbf{50.00} & 22.22 & \underline{62.99} & \textbf{45.46} \\
CausalLM & 66.90 & 25.40 & 37.60 & 24.60 & 34.00 & \textbf{33.33} & 45.70 & 38.22 \\
\midrule
\rowcolor{barGray} \multicolumn{9}{c}{\textit{Multi-turn routers}} \\
Router-R1 & 35.64 & 24.60 & 17.28 & 14.75 & 18.00 & 22.22 & 23.62 & 22.30 \\
$k$NN-MultiRound & 13.99 & 24.70 & 18.32 & 16.39 & 30.00 & 25.93 & 33.07 & 23.20 \\
LLM-MultiRound & 12.98 & 24.60 & 17.44 & 14.29 & 31.03 & 25.93 & 30.33 & 22.37 \\
\bottomrule
\end{tabular}%
}
\vspace{-10pt}
% \parbox{\linewidth}{\footnotesize \(\dagger\) All multi-turn routers use Qwen2.5-3B-Instruct as their base model.}
\end{table}

Table~\ref{tab:main} and Table~\ref{tab:personalized} report the results under the performance-only setting. Based on these results, we have the following key observations:
% {\renewcommand{\thefootnote}{$\dagger$}\footnotetext{All multi-turn routers use Qwen2.5-3B-Instruct as their base model.}}


\noindent\textbf{No single router dominates across all tasks}: The winner router varies across tasks. For example, RouterDC performs best on Generic LLM Tasks and SVMRouter on LoCoMo. Though GraphRouter attains the best average on xRouteBench, it did not consistently outperform other routers in all tasks.
% as the data-rich Generic LLM Tasks track favors the expressive RouterDC and GraphRouter, the memory sets are led by the simple embedding-based SVMRouter and $k$NNRouter, and the query-agnostic EloRouter leads the three small non-text sets, where verbalized images and series drift from the sentence-embedding distribution and a global quality rating is more robust than unreliable per-query features.

% \noindent\textbf{Multi-round routing not necessarily outperforms single-round routing}:
% Across all benchmarks, allowing the router to take multiple rounds of routing and aggregation brings no consistent gain over a single routing decision. The value of extra rounds is highly query-dependent: for many questions a single well-chosen route may be already sufficient. When a single-turn routing is sufficient to solve the problem, additional rounds may yield little benefits and instead introduce redundant informatino with unnecessary computational overhead. For the multi-turn routers, better assessing whether the information is sufficient and stopping the routing process is an important direction for improvement.
\noindent\textbf{Multi-turn routing does not consistently outperform single-turn routing}:
Across all benchmarks, multiple rounds of routing and aggregation provide no consistent gain over a single routing decision. For many queries, one well-chosen route is sufficient, whereas additional rounds introduce redundant information and computational overhead. Multi-turn routers also rely on a base model (Qwen2.5-3B-Instruct) to decompose queries and aggregate responses, making their performance sensitive to the capabilities of this model. These results highlight the need for better sufficiency estimation, early stopping, and more effective decomposition and aggregation.



\begin{wraptable}{r}{0.52\textwidth}
\vspace{-12pt}
\centering
\caption{\textbf{Performance comparison on the personalized track.} Top two results are highlighted in \textbf{bold} and \underline{underline}.}
\label{tab:personalized}
\small
\setlength{\tabcolsep}{4pt}
\begin{tabular}{lc|lc}
\toprule
\textbf{Router} & \textbf{Acc.} & \textbf{Router} & \textbf{Acc.} \\
\midrule
GMTRouter & \textbf{68.78} & RouterDC & 56.44 \\
PersonalizedRouter & \underline{67.86} & MFRouter & 54.39 \\
EloRouter & 66.40 & MLPRouter & 52.93 \\
GraphRouter & 65.23 & $k$NNRouter & 51.76 \\
SVMRouter & 65.08 & CausalLM & 46.78 \\
Largest-LLM & 58.05 & Router-R1 & 45.46 \\
Hybrid LLM & 57.91 & Smallest-LLM & 42.53 \\
\bottomrule
\end{tabular}
\vspace{-10pt}
\end{wraptable}

\noindent\textbf{Conditioning on user context helps, but how it is modeled matters}:
Table~\ref{tab:personalized} reports persona-judge accuracy on the personalized track. GMTRouter achieves the highest accuracy of 68.78, outperforming PersonalizedRouter (67.86) and the best user-agnostic router, EloRouter (66.40). The strong performance of both personalized methods confirms the benefit of conditioning routing decisions on user context, while GMTRouter's additional 0.92-point gain over PersonalizedRouter shows that the way user context is encoded and integrated remains important.



\subsection{Performance--Cost Trade-offs}
\label{sec:tradeoff}

\begin{figure}[t]
\centering
\includegraphics[width=\linewidth]{figs/rank_heatmap_main.pdf}
\caption{\textbf{Router rankings across the Generic LLM Tasks, memory, vision, and time-series tracks as the cost weight $\beta$ increases.} Each cell gives a router's rank under the weighted performance--cost objective, with smaller rank values indicating better performance.}
\label{fig:rank-beta}
\vspace{-10pt}
\end{figure}

\noindent\textbf{No single router is best under every performance--cost trade-off.} Figure~\ref{fig:rank-beta} ranks the routers by reward within each category as the cost weight $\beta$ grows, and the rankings shift dramatically along the sweep. RouterDC tops Generic LLM Tasks when only quality matters, yet falls to tenth of eleven under the most cost-sensitive setting; EloRouter leads Vision and TimeSeries at $\beta = 0$ but drops out of the lead once cost enters the objective. However, weakness at one operating point does not imply weakness at another, as MLPRouter sits near the bottom of Vision under the quality-first setting yet becomes the best choice there for every $\beta \geq 0.4$. Therefore, it is practical to choose the router that matches the performance--cost requirements of the deployment at hand.


\begin{wrapfigure}{r}{0.5\textwidth}
\centering
\vspace{-15pt}
\includegraphics[width=1\linewidth]{figs/perf_vs_price_all.pdf}
\vspace{-10pt}
\caption{\textbf{Performance--cost trade-offs of routers averaged across the xRouteBench tracks.} Each point represents an operating setting with a different cost weight $\beta$, where higher performance and lower per-query inference cost are preferred.}
\label{fig:perf-price-all}
\vspace{-30pt}
\end{wrapfigure}

\noindent\textbf{Increasing inference cost leads to improved performance.} Figure~\ref{fig:perf-price-all} presents the trade-off between performance and inference cost. For most routers, performance and cost exhibit a clear positive correlation, with the operating points rising from the low-cost to the high-cost end. This is because increasing the inference budget unlocks more powerful and expensive models, which perform better. Meanwhile, always calling the largest model incurs the highest cost yet delivers only mediocre performance and is dominated by the learned routes, since many queries that the largest model fails are solved by smaller and cheaper ones. This confirms that no single model covers all queries, which is exactly the headroom that routing exploits.


\subsection{Routing in Deployment: Real Users and Multi-Agent Systems}
\label{sec:extension-study}
\begin{wraptable}{r}{0.52\textwidth}
\vspace{-20pt}
\centering
\caption{\textbf{Router performance on held-out real-user sessions collected through the Slack deployment.} Accuracy measures how often each router's model selection agrees with the users' pairwise preferences.}
\label{tab:personalized_human}
\small
\setlength{\tabcolsep}{4pt}
\begin{tabular}{lc|lc}
\toprule
\textbf{Router} & \textbf{Acc.} & \textbf{Router} & \textbf{Acc.} \\
\midrule
PersonalizedRouter & \textbf{83.05} & RouterDC & 65.25 \\
EloRouter & \underline{82.20} & $k$NNRouter & 60.17 \\
MLPRouter & 78.81 & $k$NN-MultiRound & 60.17 \\
SVMRouter & 77.12 & Smallest-LLM & 55.08 \\
Hybrid LLM & 73.73 & MFRouter & 51.69 \\
GMTRouter & 70.70 & Largest-LLM & 41.53 \\
GraphRouter & 67.17 & CausalLM & 27.97 \\
\bottomrule
\end{tabular}
\vspace{-20pt}
\end{wraptable}

The deployment layer of \method\ carries routers beyond static benchmarks. We study two settings it enables: routing for real users served through OpenClaw and routing inside multi-agent systems.

\noindent\textbf{Routing for real users.}
Using the OpenClaw server of \method, we deploy the routing stack behind Slack and collect live preference feedback. 15 users contribute 40 sessions of 1 to 12 turns, totaling 234 pairwise records. For each query, two models are sampled from a pool of ten candidates, their answers are shown in randomized positions, and the user marks one as better or declares a tie. We split by session into 32 training sessions and 8 test sessions, train every router on the human training split, and score how often its selection matches the human preference on held-out sessions. Table~\ref{tab:personalized_human} shows that PersonalizedRouter leads at 83.05, while the fine-tuned CausalLM router ranks last. The simulated ranking does not fully transfer, as GMTRouter, the winner under the persona judge, drops to sixth on real users, showing that it matters to validate personalized routers against real feedback.

\noindent\textbf{Routing inside multi-agent systems.}
A multi-agent system (MAS) is conventionally instantiated with a single base model shared by every agent. We instead treat model choice as a per-agent decision, where a router receives the prompt of each functional node and selects the most suitable LLM for that call. Since node prompts differ substantially, covering planning, execution, and verification instructions, this setting stress-tests how well a router trained on ordinary queries generalizes. Following MultiAgentBench \citep{zhu2025multiagentbench} and GraphPlanner \citep{feng2026graphplanner}, we instantiate five coordination topologies, namely Star, Tree, Graph, Chain, and Plan-Exec-Sum, illustrated in Figure~\ref{fig:multiagent} with node roles detailed in Appendix~\ref{app:mas}. The learning-based routers are trained on the Generic LLM Tasks training split, each test query is then fed through the full MAS, and the final MAS answer is scored by the task metric. Table~\ref{tab:mas} shows that routing every node pays off, as six of the seven learned routers beat always selecting the largest model on average, with MFRouter attaining the best average of 76.48 against 71.48.

\begin{figure}[t]
\centering
\includegraphics[width=\linewidth]{figs/multiagent.pdf}
\caption{\textbf{Representative multi-agent system architectures and coordination topologies:} (a) star-based centralized coordination, (b) hierarchical tree-based delegation, (c) graph-based peer interaction, (d) sequential chain collaboration, and (e) planner–executor–summarizer workflow.}
\label{fig:multiagent}
% \vspace{-10pt}
\end{figure}

\begin{table}[t]
\centering
\caption{\textbf{Router performance on the Generic LLM Tasks test split when each node in a multi-agent system is routed independently.} Results are reported across five coordination topologies, with the final column showing the average performance.}
\label{tab:mas}
\small
\setlength{\tabcolsep}{6pt}
\begin{tabular}{lcccccc}
\toprule
\textbf{Router} & \textbf{Star} & \textbf{Tree} & \textbf{Graph} & \textbf{Chain} & \textbf{Plan-Exec-Sum} & \textbf{Avg} \\
\midrule
% Random & 73.20 & 75.20 & 78.20 & 70.00 & 73.20 & 73.96 \\
Largest-LLM & 69.00 & 67.00 & 77.20 & 69.00 & \underline{75.20} & 71.48 \\
\midrule
$k$NNRouter & 74.80 & \underline{78.60} & 78.60 & 76.60 & 71.80 & 76.08 \\
SVMRouter & \underline{76.20} & 75.60 & \underline{80.00} & 74.40 & \underline{75.20} & \underline{76.28} \\
MLPRouter & 75.40 & 76.60 & 76.80 & \underline{78.00} & 71.40 & 75.64 \\
MFRouter & 75.40 & 74.20 & \textbf{81.00} & \textbf{78.60} & 73.20 & \textbf{76.48} \\
EloRouter & 73.80 & 72.40 & 78.60 & 76.60 & \underline{75.20} & 75.32 \\
GraphRouter & 68.20 & 70.80 & 66.20 & 72.00 & 69.00 & 69.24 \\
RouterDC & \textbf{77.60} & \textbf{79.60} & 74.20 & 72.00 & \textbf{76.20} & 75.92 \\
\bottomrule
\end{tabular}
\vspace{-10pt}
\end{table}